
\documentclass[11pt]{amsart}

\usepackage{AFBoix2015}

\begin{document}

\author[A.\,F.\,Boix]{Alberto F.\,Boix}
\address{Department of Mathematics, Universitat Polit\`ecnica de Catalunya BarcelonaTech, Av. Eduard Maristany 16, 08019, Barcelona, Spain.}
\email{alberto.fernandez.boix@upc.edu}

\author[S. Zarzuela]{Santiago Zarzuela}
\address{Departament de Matem\`atiques i Inform\`atica, Universitat de Barcelona, Gran Via de les Corts Catalanes 585, 08007, Barcelona, SPAIN}
\email{szarzuela@ub.edu}

\title[On deficiency modules]{Regularity of deficiency modules through spectral sequences}

\keywords{Castelnuovo--Mumford regularity, Deficiency modules}

\subjclass[2020]{Primary 13A02, 13D07, 13D45}

\begin{abstract}
The main goal of this paper is to obtain upper bounds for the regularity of graded deficiency modules in the spirit of the one obtained by Kumini--Murai in the monomial case building upon the spectral sequence formalism developed by \`Alvarez Montaner, Boix and Zarzuela. This spectral sequence formalism allows us not only to recover Kumini--Murai's upper bound for monomial ideals, but also to extend it for other types of rings, which include toric face rings and some binomial edge rings, producing to the best of our knowledge new upper bounds for the regularity of graded deficiency modules of this type of rings.
\end{abstract}

\maketitle

\section*{Introduction}
Let $\mathbb{K}$ be a field, let $R=\mathbb{K}[x_1,\ldots, x_n]$ be a polynomial ring in $n$ variables over $\mathbb{K},$ that we assume it is standardly $\mathbb{Z}$--graded (this just means that elements of $\mathbb{K}$ have degree zero, and variables have degree one). In this setting, given a finitely generated, graded $R$--module $M$ we consider Schenzel's graded deficiency modules
\[
K^j (M):=\hbox{}^* \Ext_R^{n-j} (M,R(-n))
\]
as a way of measuring how far $M$ is of being a Cohen--Macaulay module. In this setting, it is of interest to understand how big can be the Castelnuovo--Mumford regularity of these modules; more precisely, we want to calculate upper bounds for this regularity.

This is a topic that has been considered in several research papers, where general upper bounds has been obtained (see, for instance, \cite{HoaHyry2006,BrodmannLinhSeiler2013,ChardinHaHoa2011}). However, most of these bounds are far from being explicit.

Now, we assume that $M=R/I$ where $I$ is a monomial ideal; in this setting, Kumini and Murai \cite{KumminiMurai2011} obtained the following nice and explicit upper bound:
\[
\reg (K^j (R/I))\leq j.
\]
The proof of this result deeply relies on the so--called \textit{Stanley filtrations} used for instance by Maclagan and Smith \cite{MaclaganSmith2005} to obtain upper bounds for the so--called multigraded regularity for toric varieties that they introduce in \cite{MaclaganSmith2004}; however, Stanley filtrations require working with multigradings, a technique which is almost exclusive for monomial ideals.

In other different direction, Ha and Hoa \cite[Proposition 2.3]{HaHoa2008} showed that if $I$ is an arbitrary graded ideal such that $R/I$ is Cohen--Macaulay, then
\[
\reg (K^{\dim (R/I)}(R/I))=\dim (R/I).
\]
For us, one of the interests for studying the regularity of deficiency modules is given by \cite{DaoMaVarbaro}; indeed, assume now that $I$ is a graded ideal inside $R,$ and let $(h_a (R/I),\ldots, h_s (R/I))$ be the $h$--vector of $R/I.$ By a celebrated result by Murai and Terai \cite[Theorem 1.4]{MuraiTeraihvectors}, it is known that, if there exists an integer $r\geq 0$ such that
\[
\reg (K^j (R/I))\leq j-r
\]
for all $j<\dim (R/I),$ then $h_j (R/I)\geq 0$ for all $j\leq r.$ This result was extended by Dao, Ma and Varbaro in \cite[Theorem 2.8]{DaoMaVarbaro} replacing $R/I$ by an arbitrary graded, finitely generated $R$--module $M.$ The interested reader can consult for instance Holmes work \cite[Corollary 5.5]{HolmesgeneralizedSerrecondition} for a setting where one can obtain this type of upper bounds for the regularity.

The main goal of this paper is to produce upper bounds in the spirit of the one obtained by Kumini--Murai in the monomial case building upon the spectral sequence formalism developed in \cite{AlvarezBoixZarzuelaspectralsequences}. This spectral sequence formalism allows us not only to recover Kumini--Murai's upper bound for monomial ideals, but also to extend it for other types of rings.

Now, we provide a more detailed overview of the contents of this paper for the convenience of the reader. After reviewing in Section \ref{section on preliminaries on spectral sequences} some preliminary results obtained in \cite{AlvarezBoixZarzuelaspectralsequences} that we plan to use along this paper, in Section \ref{section on ultrametric functions} we introduce what we call ultrametric functions, that can be regarded as a generalization of the notion of regularity. We do so mainly in order to illustrate that the formalism used throughout this paper can be applied not only for the regularity, but also for other functions that behave, roughly speaking, like regularity on short exact sequences.

Section \ref{section of main results} contains the main results of this paper; here, we want to single out Theorem \ref{the expected bound when P is a C.M. poset}, where we establish our main upper bound for the regularity of graded deficiency modules. Theorem \ref{the expected bound when P is a C.M. poset} recovers and extends \cite{KumminiMurai2011} beyond the monomial case.

In the remainder sections we specialize Theorem \ref{the expected bound when P is a C.M. poset} to some concrete settings; more precisely, in Section \ref{section: the case of toric face rings} we specialize to the case of toric face rings, obtaining, to the best of our knowledge, a new upper bound for the regularity of deficiency modules of a toric face ring (see Theorem \ref{deficiency in case of toric face rings}). On the other hand, in Section \ref{section: the case of binomial edge ideals} we also obtain similar upper bounds for some binomial edge rings (see Theorem \ref{deficiency in case of binomial edge ideals}). Finally, in Section \ref{section of examples} we illustrate our main results by means of examples.

%The same can be said about the upper bound for the regularity of deficiency modules in the case of binomial edge ideals obtained in Section \ref{section: the case of binomial edge ideals}.

\section{Preliminaries on some spectral sequences}\label{section on preliminaries on spectral sequences}

\subsection{Construction of homological spectral sequences}
The purpose of this section is to collect some results borrowed from \cite{AlvarezBoixZarzuelaspectralsequences} that we plan to use along this paper.
The setup we need for the appropriate  spectral sequence is the following:



\begin{cons}\thlabel{otra construccion mas de lo mismo}
Let $\mathbb{K}$ be a field, let $A=\mathbb{K}[x_1,\ldots ,x_n]$ be the polynomial ring in $n$ variables over $\mathbb{K}$ standardly $\mathbb{Z}$--graded, and let $\mathcal{A}$ be the category of fintely generated graded $A$--modules. Moreover, let $\xymatrix@1{\cA\ar[r]^-{T}& \cA}$ be a contravariant, left exact functor, let $P$ be any finite poset, let $\Inv (P,\cA)$ be the category of inverse systems on $P$ valued on $\cA$, and let $\Dir (P,\cA)$ be the category of direct systems on $P$ valued on $\cA$.
Carrying over $T$, we produce a new functor (namely, $\mathcal{T}$) in the following manner:
\begin{align*}
& \xymatrix{\Inv (P,\cA)\ar[r]^-{\mathcal{T}}& \Dir (P,\cA)}\\ & G=(G_p)_{p\in P}\longmapsto
\left(T\left(G_p\right)\right)_{p\in P}.
\end{align*}
Moreover, we also assume that $T$ commutes with finite direct sums and that $T(A)=Z$ for some $A$-module $Z$;
in particular, we have that $\colim_{p\in P} \mathcal{T} \left(A_{\leq q}\right)=Z$, where $q\in P$ (notice that $Z$
only depends on $T$, but not on $q$). Recall that by $A_{\leq q}$ we mean the following inverse system.
\[
(A_{\leq p})_q=\begin{cases}
    A,\text{ if }q\leq p,\\
    0,\text{ otherwise.}
\end{cases}
\]
On the other hand, given $q\in P$ and $M$ an object of $\mathcal{A}$, we denote by $M_q$ the following direct system: for any $p\in P,$
\[
(M_q)_p=\begin{cases}
M,\text{ if }p=q,\\
0,\text{ otherwise.}
\end{cases}
\]
Finally, we define the cohomological dimension of the functor $T$ as follows:
\[
\cdim(T):=\max\{i\in\mathbb{N}:\ \mathbb{R}^i T(X)\neq 0\text{ for some graded $A$--module }X\}.
\]
\end{cons}

Next result can be regarded as a corrected version of the statement of \cite[Theorem 6.5 and Corollary 6.6]{AlvarezBoixZarzuelaspectralsequences}. Indeed, we want to single out that the  difference between the statement of \cite[Theorem 6.5 and Corollary 6.6] {AlvarezBoixZarzuelaspectralsequences} and the one presented here is part (b) of our assumptions. We reproduce the proof for the convenience of the reader. 

\begin{teo}\thlabel{sucesion espectral y colapso todo en uno: segunda parte}
Preserving the notations and assumptions established along \thref{otra construccion mas de lo mismo}, we further assume that the inverse
system $V$ is acyclic with respect to the limit on $P$, and that there is a natural equivalence of functors
\[
\colim_{p\in P}\circ\mathcal{T}\cong T\circ\lim_{p\in P}.
\]
Then, the following assertions hold.

\begin{enumerate}[(a)]

\item There is a spectral sequence of the form
\[
E_2^{-i,j}=\mathbb{L}_i \colim_{p\in P} \mathbb{R}^j \mathcal{T} (V)\Longrightarrow \mathbb{R}^{j-i} T\left(\lim_{p\in P} V\right).
\]

\item Assume, in addition, that for any $p\in P$, $\R^j T \left(V_p\right)=0$ up to a single value of $j$ (namely, $h_p$), and that there is a canonical isomorphism of graded $A$--modules
\[
\mathbb{R}^j \mathcal{T} (V)\cong\bigoplus_{j=h_q} \left(\mathbb{R}^{h_q} T (V_q)\right)_q.
\]
Then, the previous spectral sequence can be expressed in the following way:
\[
E_2^{-i,j}=\bigoplus_{j=h_q} \R^{h_q} T \left(V_q\right)^{\oplus m_{i,q}}\xymatrix{ \ar@{=>}[r]_-{i}& }\R^{j-i}T \left(\lim_{p\in P}V_p\right),
\]
where $m_{i,q}:=\dim_{\K} \left(\widetilde{H}_{i-1} \left(\left(q,1_{\widehat{P}}\right);\K\right)\right)$. Moreover,
this spectral sequence degenerates at the $E_2$-sheet.

Therefore, by \cite[5.2.5]{Weibel1994}, for each
$0\leq r\leq\cdim (T)$ there is an increasing, finite filtration $\{H_k^r\}$ of $\R^r T(\lim_{p\in P}V_p)$ by graded
$A$-modules such that, for any $k\geq 0$,
\[
H_k^r/H_{k-1}^r\cong E_2^{-k,r+k}=\bigoplus_{\{q\in P\ \mid\ r+k=h_q\}} \R^{h_q} T \left(V_q\right)^{\oplus m_{k,q}},
\]
where we follow the convention that $H_{-1}^r =0$.

\end{enumerate}
\end{teo}

\begin{proof}
First of all, we prove part (i).  Let $\xymatrix@1{\ldots\ar[r]& F_2\ar[r]& F_1\ar[r]& F_0\ar[r]& V\ar[r]&  0}$ be a projective resolution of $V$ in the category $\Inv (P,\cA)$ such that, for any $i\geq 0$, $F_i$ can be expressed as a finite direct sum of inverse systems of the form $(A(n))_{\leq q}$ for some $q\in P$ and some integer $n\in\mathbb{Z}$ (see \cite[Theorem 2.17]{AlvarezBoixZarzuelaspectralsequences}). Now, by applying the contravariant functor $\mathcal{T}$ we obtain the following cochain complex in the category $\Dir (P,\cA)$:
\[
\xymatrix{0\ar[r]& \mathcal{T} (V)\ar[r]& \mathcal{T}(F_0)\ar[r]& \mathcal{T} (F_1)\ar[r]& \ldots}
\]
Now, by applying the homological Roos complex \cite[\S 3.1]{AlvarezBoixZarzuelaspectralsequences} for computing the left derived functors of
the colimit, we end up with the following bicomplex:
\[
\xymatrix{\hbox{}& \vdots& \vdots& \hbox{}\\
\ldots\ar[r]& \roos^{-1}(\mathcal{T} (F_1))\ar[u]\ar[r]& \roos^0(\mathcal{T} (F_1))\ar[u]\ar[r]& 0\\
\ldots\ar[r]& \roos^{-1}(\mathcal{T} (F_0))\ar[u]\ar[r]& \roos^0(\mathcal{T} (F_0))\ar[u]\ar[r]& 0\\
\hbox{}& 0\ar[u]& 0\ar[u]& \hbox{}}
\]
In other words, we use the homological Roos complex in order to produce the bicomplex
\[
\roos^{-i} \left(\mathcal{T}\left(F_j\right)\right)=\roos_i
\left(\mathcal{T}\left(F_j\right)\right);
\] the reader will easily note
that we put a minus in the index $i$ because we want to work with a
bicomplex located in the second quadrant.

Moreover, we have to stress that the horizontal differentials are the
ones of the homological Roos complex and the vertical ones are the
induced by the projective resolution of $V$; so, the bicomplex
$\roos^{-i} \left(\mathcal{T}\left(F_j\right)\right)$ produces two spectral
sequences; namely, the ones provided respectively by the first and
the second filtration of the previous bicomplex. In this way, the
first thing one should ensure is that both spectral sequences
converge and calculate their common abutment.

First of all, we want to check that our spectral sequences converge;
indeed, since $\roos^{-i} (\mathcal{T} (F_j))=0$ for all $i\gg 0$, we have that the bicomplex $\roos^{-i} \left(\mathcal{T}\left(F_j\right)\right)$ has just a finite number of
columns, which implies the convergence of both spectral sequences.

Second, notice that the $E_2$-page of one of these spectral
sequences is obtained by firstly computing the homology of the
rows and then computing the cohomology of the columns; indeed, \cite[Lemma 6.2]{AlvarezBoixZarzuelaspectralsequences} guarantees that all the
rows of the bicomplex $\roos^{-i} \left(\mathcal{T}\left(F_j\right)\right)$ are exact up to the $0$th spot. More precisely, in this $0$th spot we obtain $(\colim_{p\in P}\circ\mathcal{T})(F_j)$ for each $j$. Now, by using the natural equivalence
\[
\colim_{p\in P}\circ\mathcal{T}\cong T\circ\lim_{p\in P}
\]
given by the assumption we have that all the rows of the bicomplex are exact up to the $0$th spot, and in this position we obtain $(T\circ\lim_{p\in P}) (F_j)$ for each $j$. Now, keeping in mind this fact, we obtain (after computing cohomology on the columns) that the abutment is equal to $\R^* (T\circ\lim_{p\in P}) (V)$.

Now, our next claim is that, for each $F_j$ as above, $\lim_{p\in P} F_j$ is an acyclic object with respect to $T$. Indeed, take the inverse system $(A(t))_{\leq q}$, where $q\in P$ and $t\in\mathbb{Z}$. In this case, it is clear that
\[
\lim_{p\in P} (A(t))_{\leq q}=A(t).
\]
Now, we come back to the general case, since each $F_j$ appearing in the resolution of $V$ is the finite direct sum of inverse systems of the form $(A(t))_{\leq q}$ and the inverse limit commutes with finite direct sums we finally obtain that $\lim_{p\in P} F_j$ can be expressed as finite direct sums of objects of the form $A(t)$, which are acyclic with respect to $T$.

Now, since $V$ is acyclic by assumption with respect to the limit on $P$, we have that
\[
\xymatrix{\ldots\ar[r]& }\lim_{p\in P}F_2\xymatrix{\ar[r]& }\lim_{p\in P}F_1\xymatrix{\ar[r]& }\lim_{p\in P}F_1\xymatrix{\ar[r]& }\lim_{p\in P}V\xymatrix{\ar[r]& 0}
\]
is a resolution of $\lim_{p\in P}V$ by $\hbox{}^*$ free objects in the category of graded $A$--modules, and therefore we finally end up, for each $i$, with the natural isomorphism
\[
(\R^i T) \left(\lim_{p\in P} V\right)\cong\R^{i} \left(T\circ\lim_{p\in P}\right) (V).
\]
Summing up, we end up with the fact that the abutment is precisely $\R^* T \left(\lim_{p\in P} V\right)$.

Finally, the other spectral sequence that we can produce is the one
obtained by firstly taking cohomology on the columns and then
calculating the cohomology of the rows; in this case, one obtains
as $E_1$-page $E_1^{-i,j}=\roos^{-i} \left(\R^j
\mathcal{T} (V)\right)$. In addition, since the boundary map of the
$E_1$-page is the one of the Roos chain complex, and this chain
complex computes the $i$th left derived functor of the colimit, its
$E_2$-page turns out to be $E_2^{-i,j}=\Lf_i \colim_{p\in P} \R^j
\mathcal{T} (V)$. Summing up, combining all the foregoing facts one obtains
the spectral sequence
\[
E_2^{-i,j}=\mathbb{L}_i \colim_{p\in P} \mathbb{R}^j \mathcal{T} (V)\Longrightarrow \mathbb{R}^{j-i} T\left(\lim_{p\in P} V\right).
\]
in the category of graded $A$-modules; this proves part (i).

Now, let us prove part (ii). By our assumption, we have a canonical isomorphism of graded $A$--modules
\[
\mathbb{R}^j \mathcal{T} (V)\cong\bigoplus_{j=h_q} \left(\mathbb{R}^{h_q} T (V_q)\right)_q.
\]
Now, fix $i\in\mathbb{N}_0$. Applying to the above isomorphism the $i$th left derived functor of the colimit over $P$, we obtain the following canonical isomorphism of graded $A$--modules:
\[
\mathbb{L}_i \colim_{p\in P} \mathbb{R}^j \mathcal{T} (V)\cong\bigoplus_{j=h_q} \mathbb{L}_i \colim_{p\in P}\left(\mathbb{R}^{h_q} T (V_q)\right)_q.
\]
Now, by \cite[Lemma 4.10]{AlvarezBoixZarzuelaspectralsequences} we know that
\[
\mathbb{L}_i \colim_{p\in P}\left(\mathbb{R}^{h_q} T (V_q)\right)_q\cong\bigoplus\widetilde{H}_{i-1} ((q,1_{\widehat{P}}); \mathbb{R}^{h_q} T (V_q)).
\]
Now, by giving the same arguments explained along the proof of \cite[Theorem 4.11]{AlvarezBoixZarzuelaspectralsequences} we conclude that there is a canonical isomorphism of graded $A$--modules
\[
\mathbb{L}_i \colim_{p\in P}\left(\mathbb{R}^{h_q} T (V_q)\right)_q\cong\bigoplus_{j=h_q} \R^{h_q} T \left(V_q\right)^{\oplus m_{i,q}},
\]
just what we finally wanted to show.
\end{proof}

%Now, consider the following Grothendieck's spectral sequence \cite[Theorem 10.50]{Rotman2009}:
%\[
%(\R^i T)\left(\R^j\lim_{p\in P}\right)(V)\xymatrix{\ar@{=>}[r]& }\R^{i+j} \left(T\circ\lim_{p\in P}\right) (V).
%\]

\section{On ultrametric functions}\label{section on ultrametric functions}

The goal of this section is to introduce the notion of the so--called ultrametric functions, mainly because the regularity of a finitely generated, graded module over a polynomial ring may be regarded as a particular example of ultrametric function. We want to mention here that this definition is not standard terminology and it is only coined for the purposes of this paper. The reader may like to compare our definition of ultrametric function with the notion of Sylvester matrix function studied in \cite{Virilirankfunctions}, and also with the approach developed in \cite[Lemma 2.1]{Santoni}.

\begin{df}\thlabel{introducing ultrametric functions}
Let $\mathcal{A}$ be an abelian category, let $G$ be a finitely generated, torsion free abelian group with $r=\operatorname{rank}(G)$, let $P(G)$ be the poset of all subsets of $G$, and let $\xymatrix@1{\mathcal{A}\ar[r]^-{\mu}& \widehat{M}}$ be a map, where
\[
\widehat{M}=\begin{cases} \mathbb{Z}\cup\{-\infty\},\text{ if }r=1,\\
P(G),\text{ if }r\geq 2.
\end{cases}
\]
We say that $\mu$ is an \textit{ultrametric function} if
\[
\mu (0)=\begin{cases} -\infty,\text{ if }r=1,\\
G,\text{ if }r\geq 2.
\end{cases}
\]
and, for any short exact sequence $\xymatrix@1{0\ar[r]& M'\ar[r]& M\ar[r]& M''\ar[r]& 0}$ in $\mathcal{A}$, we have
\[
\mu (M)\begin{cases} \leq\max\{\mu (M'),\ \mu (M'')\},\text{ if }r=1,\\
\supseteq \mu(M')\cap\mu (M''),\text{ if }r\geq 2.\end{cases}
\]
\end{df}
Here, we present some examples of ultrametric functions.

\begin{ex}\thlabel{examples of ultrametric functions}
The goal here is to introduce several examples of ultrametric functions that appear quite naturally in the realm of Commutative Algebra and Algebraic Geometry.

\begin{enumerate}[(i)]

\item Let
\[
R=\bigoplus_{n\geq 0} R_n
\]
be a positively graded ring such that $R=R_0 [R_1],$ and let $\mathcal{A}$ be the category of finitely generated, graded $R$--modules. For each integer $l\geq 0,$ one has, as consequence of \cite[16.2.15 (ii)]{BroSha}, that the Castelnuovo--Mumford regularity $\reg^l$ at and above level $l$ is an ultrametric function of $\mathcal{A}.$

\item Under the same assumptions of part (i), let $\mathcal{X}\subset\mathbb{N}$ be a set of indices. We have, by \cite[Lemma 2.3 (i)]{Cavigliatensorproducts} that the regularity $\reg^{\mathcal{X}}$ with respect to $\mathcal{X}$ defined by Caviglia in \cite[Definition 2.1]{Cavigliatensorproducts} is also an ultrametric function.

\item Let now $\mathcal{C}$ be the category of finitely generated modules over a Noetherian ring $R,$ and let $I\subset R$ be any ideal. Then, by \cite[Proposition 1.2.9]{BrunsHerzog1993} we know that $\mu=-\grade_I$ is also an ultrametric function.

\item Let now $\mathcal{B}$ be the category of finitely generated modules over a local ring $(R,\mathfrak{m})$ having finite projective dimension. Then, the projective dimension $\projdim$ is also an ultrametric function on $\mathcal{B}.$

\item More generally, let now $\mathcal{G}$ be the category of finitely generated modules over a commutative Noetherian ring having finite Gorenstein dimension \cite[(1.2.3)]{Gdimensionbook}. Then, the Gorenstein dimension $\Gdim$ is an ultrametric function on $\mathcal{G}$ because of \cite[(1.2.9)]{Gdimensionbook}. We recall that, for finitely generated modules with finite projective dimension, projective and Gorenstein dimension coincide, as it is shown in \cite[(1.2.10)]{Gdimensionbook}.

\item Let now $G$ be a finitely generated, torsion free abelian group, let $R$ be a polynomial ring over a field graded by $G,$ and let $\mathcal{C}$ be the category of finitely generated, graded $R$--modules. In this case, the Castelnuovo--Mumford regularity as defined by Maclagan and Smith in \cite{MaclaganSmith2004} is also an ultrametric function because of \cite[Lemma 7.1]{MaclaganSmith2004}.

\item Let $\mathcal{C}$ be the category of finitely generated $R$--modules, where $R$ is a commutative Noetherian ring, and let $I$ be an ideal of $R.$ The cohomological dimension $\cdim(I,-)$ with respect to $I$ is also an ultrametric function, as proved for instance by Divaani--Aazar, Naghipour and Tousi in \cite[Corollary 2.3 (i)]{DivaaniAazarNaghipourTousi2002}.

\item Let $R$ be a commutative Noetherian local ring of prime characteristic $p,$ let $F$ be the Frobenius map on $R,$ and denote by $\mathcal{C}$ the category of left $R[\Theta; F]$--modules. Denote by $\rho$ the $F$--depth (respectively, generalized $F$--depth) function of an object of $\mathcal{C}$, which have been studied for instance in the work by Maddox and Miller \cite{MaddoxMiller}. Then, thanks to \cite[Theorem 3.5]{MaddoxMiller}, one has that $\mu=-\rho$ is also an ultrametric function.

\end{enumerate}

\end{ex}
One of the main reasons for considering ultrametric functions is their behaviour with respect to filtrations; namely:

\begin{lm}\thlabel{ultrametric functions and filtrations}
Let $\mathcal{A}$ be an abelian category, let $\mu$ be an ultrametric function on $\mathcal{A},$ set $r:=\rank (G)$ as in \thref{introducing ultrametric functions}, and let $0=F_{-1}\subset F_0\subset F_1\subset\ldots\subset F_t=F$ be a filtration in $\mathcal{A}.$ Then, one has that, if $r=1,$ then
\[
\mu (F)\leq\max_{0\leq k\leq t}\left\{\mu\left(F_k/F_{k-1}\right)\right\}.
\]
On the other hand, if $r\geq 2$ then
\[
\mu (F)\supseteq\bigcap_{k=0}^t \mu\left(F_k/F_{k-1}\right).
\]
\end{lm}

\section{Main results}\label{section of main results}
As the title says, the main purpose of this section is to present the main results of this paper. Before that, we establish the following consequence of Theorem \ref{sucesion espectral y colapso todo en uno: segunda parte}.

\begin{teo}\thlabel{spectral sequence for the deficiency modules}
Let $\mathbb{K}$ be a field, let $A=\mathbb{K}[x_1,\ldots ,x_n]$ be the polynomial ring in $n$ variables over $\mathbb{K}$ standardly $\mathbb{Z}$--graded, and let $\mathcal{A}$ be the category of finitely generated, graded $A$--modules. Moreover, set $\omega_R:=A(-n),$ let $V$ be an inverse system of $\mathcal{A}$ indexed by a poset $P$ which is acyclic with respect to the limit on $P$. Finally, given $T=\hbox{}^*\Hom_A (-,A(-n))$, we denote by $\hbox{}^*\mathcal{H}om (-,\lvert A(-n)\rvert)$ the corresponding contravariant functor from inverse systems to direct systems of $P$ (see \thref{otra construccion mas de lo mismo}). Then, there exists a spectral sequence
\[
E_2^{-i,j}=\mathbb{L}_i \colim_{p\in P} \hbox{}^*\mathcal{E}xt_R^j (V_p,\lvert \omega_R\rvert)\xymatrix{\hbox{}\ar@{=>}[r]& }\hbox{}^*\Ext_R^{j-i} \left(\lim_{p\in P}V_p,\omega_R\right).
\]
\end{teo}

\begin{disc}\thlabel{the inequality for any ultrametric function}
Let $I$ be a homogeneous ideal of $A$, let $P$ be the poset given by the all the possible sums of the primary components of $I$ ordered by reverse inclusion, and let $A/[*]$ be the inverse system given, for each $p\in P$, by $A/I_p$, where the structural maps are either zero or the projections $A/I_p\longrightarrow A/I_q$ when $q<p$. Assume that the below statements hold.

\begin{enumerate}[(i)]

\item $P$ is a subset of a distributive lattice of ideals of $A$. Notice that this assumption implies \cite[Proposition 5.1]{AlvarezBoixZarzuelaspectralsequences} that the inverse system $A/[*]$ is acyclic with respect to the limit on $P$.

\item For any $p\in P,$ $A/I_p$ is a Cohen--Macaulay ring. Observe that this assumption implies that $\hbox{}^*\Ext_A^k (A/I_p,A(-n))$ is zero up to a single value of $k$.

\item If $p<q$ then $\height(I_q)<\height(I_p)$. This assumption ensures that there is a canonical isomorphism of graded $A$--modules
\[
\hbox{}^*\mathcal{E}xt_A^j(A/[*],\lvert A(-n)\rvert)\cong\bigoplus_{j=h_q} \left(\hbox{}^*\Ext_A^{h_q} \left(A/I_q,A(-n)\right)\right)_q.
\]

\end{enumerate}
In this way, the spectral sequence obtained in \thref{spectral sequence for the deficiency modules} deduced from \thref{sucesion espectral y colapso todo en uno: segunda parte} boils down to the below one:
\[
E_2^{-i,j}=\bigoplus_{j=h_q} \hbox{}^*\Ext_A^{h_q} \left(V_q,A(-n)\right)^{\oplus m_{i,q}}\xymatrix{ \ar@{=>}[r]_-{i}& }\hbox{}^*\Ext_A^{j-i} \left(\lim_{p\in P}V_p,A(-n)\right),
\]
where $m_{i,q}:=\dim_{\K} \left(\widetilde{H}_{i-1} \left(\left(q,1_{\widehat{P}}\right);\K\right)\right)$.

Now, given a finitely generated, graded $A$--module, its $j$th deficiency graded module is given by
\[
K^j (M):=\hbox{}^*\Ext_A^{n-j} (M,A(-n)).
\]
Then, appealing again to \thref{sucesion espectral y colapso todo en uno: segunda parte}, we have that for any integer $j\geq 0$ there exists a filtration $\{H_k^{n-j}\}$ of
\[
K^j \left(\lim_{p\in P} A/I_p\right)
\]
such that, for each $k\geq 0,$
\[
H_k^{n-j}/H_{k-1}^{n-j}\cong\bigoplus_{\substack{p\in P\\ n-j+k=h_p}} K^{\dim (A/I_p)} (A/I_p)^{\oplus m_{n-j,p}},
\]
where we follow as usual the convention $H_{-1}^{n-j}=0.$ This is equivalent to say, since $n=h_p+\dim (A/I_p),$ that
\[
H_k^{n-j}/H_{k-1}^{n-j}\cong\bigoplus_{\substack{p\in P\\ j-k=\dim (A/I_p)}} K^{\dim (A/I_p)} (A/I_p)^{\oplus m_{n-j,p}},
\]
where
\[
m_{n-j,p}:=\dim_{\mathbb{K}}\left(\widetilde{H}_{j-\dim(A/I_p)-1}((p,1_{\widehat{P}});\mathbb{K})\right).
\]
In this way, given any ultrametric function $\mu$ on the category of finitely generated, $\mathbb{Z}$--graded $A$--modules we have, according to \thref{ultrametric functions and filtrations}, that
\begin{equation}\label{general inequality}
\mu\left(K^j \left(\lim_{p\in P} A/I_p\right)\right)\leq\max_{p\in S_j}\mu\left(K^{\dim (A/I_p)} (A/I_p)\right),
\end{equation}
where $S_j:=\{p\in P:\ j\geq\dim (A/I_p),\ m_{n-j,p}\neq 0\}$.
\end{disc}

Next statement is the main result of this paper, which recovers and extends Kumini--Murai's upper bound for the regularity of deficiency modules of monomial ideals.

\begin{teo}\thlabel{the expected bound when P is a C.M. poset}
Preserving the assumptions and notations fixed along \thref{the inequality for any ultrametric function}, we have that
\[
\reg\left(K^j\left(\lim_{p\in P} A/I_p\right)\right)\leq\max_{p\in S_j} (\dim (A/I_p))\leq j.
\]
\end{teo}

\begin{proof}
First of all, without loss of generality we can assume that $\depth(A/I)\leq j\leq\dim (A/I),$ otherwise $K^j (A/I)=0$ and the result is obviously true. As we observed in Example \ref{examples of ultrametric functions}, the regularity is an ultrametric function; therefore, according to \thref{the inequality for any ultrametric function} we have that
\[
\reg\left(K^j\left(\lim_{p\in P} A/I_p\right)\right)\leq\max_{p\in S_j}\reg(K^{\dim (A/I_p)}(A/I_p)).
\]
Moreover, since the poset $P$ satisfies, for any $p\in P,$ that $A/I_p$ is a Cohen--Macaulay ring and therefore, by applying \cite[Proposition 2.3]{HaHoa2008}, we obtain that
\[
\reg(K^{\dim (A/I_p)}(A/I_p))=\dim (A/I_p),
\]
just what we wanted to show.
\end{proof}

\begin{rk}
The reader will easily note that, given $j\geq 0$ an integer,
\[
\{p\in P:\ j=\dim (A/I_p),\ m_{n-j,p}\neq 0\}=\{p\in P:\ j=\dim (A/I_p)\}
\]
if and only if $(p,1_{\widehat{P}})=\emptyset;$ indeed, if $j=\dim (A/I_p)$ then
\[
\widetilde{H}_{j-\dim (A/I_p)-1}((p,1_{\widehat{P}});\mathbb{K})=\widetilde{H}_{-1}((p,1_{\widehat{P}});\mathbb{K}),
\]
and this homology group is not zero if and only if $(p,1_{\widehat{P}})=\emptyset,$ as claimed. From here and from \thref{the expected bound when P is a C.M. poset} we can deduce some interesting consequences.

\begin{enumerate}[(i)]

\item On the one hand, given $j\geq 0$, assume that there exists $p\in P$ such that $I_p$ is minimal and $\dim (A/I_p)=j$. This single fact implies that $K^j (A/I_p)\neq 0$.

\item On the other hand, in case $\reg (K^j (A/I))=j$, then we can guarantee that there exists $p\in P$ such that $\dim (A/I_p)=j$ and $m_{n-j,p}\neq 0$.

\end{enumerate}
\end{rk}

\begin{rk}
Let $\mathbb{K}$ be a field, let $A=\mathbb{K}[x_1,\ldots,x_n],$ graded in the standard way, and let $M$ be a finitely generated, graded $A$--module of dimension $d$. Borrowing terminology from \cite[Definition 2.3]{DaoMaVarbaro}, we say that $M$ satisfies the $r$--th Murai--Terai's condition $(MT_r)$ provided $\reg (K^i (M))\leq i-r$ for any $0\leq i\leq d-1$.

The reader will easily note that, under the assumptions of \thref{the inequality for any ultrametric function}, \thref{the expected bound when P is a C.M. poset} shows in particular that $A/I$ always satisfies $(MT_0)$. It is also shown in \cite[Propositions 3.1 and 4.1]{DaoMaVarbaro} that $A/I$ also satisfies $(MT_0)$ if either $A/I$ defines a Du Bois singularity in characteristic zero, or an $F$--pure singularity in prime characteristic.
\end{rk}

\section{The case of toric face rings}\label{section: the case of toric face rings}

The purpose of this section is to specialize \thref{the expected bound when P is a C.M. poset} in the case of toric face rings. Indeed, the following upper bound for the regularity of deficiency modules of certain toric face rings is, to the best of our knowledge, completely new. Once again, the following statement is an immediate consequence of \thref{the expected bound when P is a C.M. poset}. For unexplained terminology concerning toric face rings, we refer the reader, for instance, to Nguyen's paper \cite{HopNguyen2012}.

\begin{teo}\thlabel{deficiency in case of toric face rings}
Let $\Sigma\subseteq\mathbb{R}^d$ be a rational pointed fan, let $\mathcal{M}_{\Sigma}$ be a Cohen--Macaulay monoidal complex supported on $\Sigma,$ let $\mathbb{K}$ be a field, and let $B:=\mathbb{K}[\mathcal{M}_{\Sigma}]$ be the corresponding toric face ring. Moreover, we write $B:=A/I,$ where $A:=\mathbb{K}[x_1,\ldots, x_n]$ and $I$ is an ideal generated by squarefree monomials and pure binomials. Then, we have that
\[
\reg(K^j(B))\leq j.
\]
\end{teo}

\begin{proof}
We only need to check that conditions (i), (ii) and (iii) established along \thref{the inequality for any ultrametric function} hold. On the one hand, condition (i) was proved by Brun, Bruns and R\"{o}mer in \cite[Proposition 2.2 and Proof]{BrunBrunsRomer2007}. On the other hand, conditions (ii) and (iii) also hold because we are assuming that our monoidal complex is Cohen--Macaulay. The proof is therefore completed.
\end{proof}

\begin{rk}
The statement of \thref{deficiency in case of toric face rings} provides, to the best of our knowledge, a new bound for the regularity of these deficiency modules. Notice that our assumptions do not necessarily imply that our toric face ring is Cohen--Macaulay, see for instance \cite[Theorem 1.1]{IchimRomer2007} obtained by Ichim and R\"{o}mer.
\end{rk}

\section{The case of binomial edge ideals}\label{section: the case of binomial edge ideals}
The purpose of this section is to specialize our main results to obtain upper bounds of the regularity of the deficiency modules of the so--called binomial edge ideals, that were independently introduced in Herzog et al. \cite{binomialedgeideals} and Ohtani \cite{Ohtanibinomialedgeideals}, we refer to the interested reader to consult the recent book \cite[Chapter 7]{HerzogHibiOhsugibinomialidealsbook} for more information concerning this specific class of ideals. For our purposes here, we are interested in the following construction, which is borrowed from work by \`Alvarez Montaner (see \cite[Definition 3.3]{AlvarezMontaner2020}).

\begin{cons}\thlabel{Alvarez Montaner poset: construction}
Let $\mathbb{K}$ be a field, let $A=\mathbb{K}[x_1,\ldots,x_n,y_1,\ldots,y_n]$ be a polynomial ring over $\mathbb{K},$ and let $I\subseteq A$ be the binomial edge ideal associated to a graph $G$ in the set of vertices $[n].$ Moreover (see \cite[Theorem 7.16]{HerzogHibiOhsugibinomialidealsbook}), let
\[
J_G=P_{D_1} (G)\cap\ldots\cap P_{D_r} (G)
\]
be the minimal prime decomposition of $J_G$ (indeed, this is possible because $J_G$ is a radical ideal, see \cite[Corollary 7.13]{HerzogHibiOhsugibinomialidealsbook}) and let $P$ be the poset given by all the possible sums of the ideals in the decomposition ordered by reverse inclusion.

In this way, we construct in an iterative way a new poset $Q:=Q_{J_G}$ as follows: it is made up by the prime ideals contained in $P,$ and by the prime ideals appearing in the posets $P_{I_q}$ of sums of ideals in the decomposition of every non prime ideal $I_q$ of $P,$ and the prime ideals we obtain repeating this procedure every time we find a non prime ideal.
\end{cons}

Keeping in mind \thref{Alvarez Montaner poset: construction}, the main result of this section can be stated as follows.

\begin{teo}\thlabel{deficiency in case of binomial edge ideals}
Let $\mathbb{K}$ be a field, let $A=\mathbb{K}[x_1,\ldots,x_n,y_1,\ldots,y_n]$ be a polynomial ring over $\mathbb{K},$ and let $I\subseteq A$ be the binomial edge ideal associated to a graph $G$ in the set of vertices $[n]$. Assume that the inverse system $(A/I_q)_{q\in Q}$ is acyclic with respect to the limit functor. Then, we have that
\[
\reg\left(K^j\left(\lim_{q\in Q}A/I_q\right)\right)\leq j.
\]
\end{teo}

\begin{proof}
We only need to check that conditions (i), (ii) and (iii) established along \thref{the inequality for any ultrametric function} hold for the poset $Q$. Indeed, first of all condition (i) follows because we have supposed that $(A/I_q)_{q\in Q}$ is acyclic with respect to the limit functor. On the other hand, parts (ii) and (iii) are also true as observed along \cite[page 338, paragraph before Discussion 3.6]{AlvarezMontaner2020}. The proof is therefore completed.
\end{proof}

\begin{rk}\label{Alvarez Montaner remark}
There is a gap in the proof of \cite[Proposition 3.7]{AlvarezMontaner2020}, and thus we cannot guarantee the acyclicity with respect to the limit functor of the inverse system $(A/I_q)_{q\in Q}$. Is for this reason that we require this assumption in the statement of Theorem \ref{deficiency in case of binomial edge ideals}.
\end{rk}

\section{Examples}\label{section of examples}
The purpose of this final section is to illustrate several of our results by means of examples. In all these examples the unjustified claims and calculations were done with Macaulay2 \cite{M2}. Our first example illustrates that, in general, our bound is sharper for Kumini--Murai's one even in the monomial case.

\begin{ex}\thlabel{sometimes S_j is empty}
Let $I=(xz,xw,yz,yw)\subset R:=\mathbb{K}[x,y,z,w].$ In this case, we know that $I=(x,y)\cap (z,w),$ and that our poset has vertices
\[
(x,y)=p_1,\ (z,w)=p_2,\ p_3=(x,y,z,w).
\]
Moreover, the deficiency modules $K^j (R/I)$ are non--zero for $j=1$ and $j=2$. Moreover, we also have $(p_1,1_{\widehat{P}})=\emptyset=(p_2,1_{\widehat{P}}),\ (p_3,1_{\widehat{P}})=\{p_1\}\cup\{p_2\}$. This implies that
\begin{align*}
& \widetilde{H}_{j-\dim (R/I_{p_1})-1}((p_1,1_{\widehat{P}});\mathbb{K})\neq 0\Longleftrightarrow j=2,\\
& \widetilde{H}_{j-\dim (R/I_{p_2})-1}((p_2,1_{\widehat{P}});\mathbb{K})\neq 0\Longleftrightarrow j=2,\\
& \widetilde{H}_{j-\dim (R/I_{p_3})-1}((p_3,1_{\widehat{P}});\mathbb{K})\neq 0\Longleftrightarrow j=1.
\end{align*}
This shows that
\begin{align*}
& S_0=\{p\in P:\ 0\geq\dim (A/I_p),\ \widetilde{H}_{0-\dim(A/I_p)-1}((p,1_{\widehat{P}});\mathbb{K})\neq 0\}=\emptyset,\\
& S_1=\{p\in P:\ 1\geq\dim (A/I_p),\ \widetilde{H}_{1-\dim(A/I_p)-1}((p,1_{\widehat{P}});\mathbb{K})\neq 0\}=\{p_3\},\\
& S_2=\{p\in P:\ 2\geq\dim (A/I_p),\ \widetilde{H}_{2-\dim(A/I_p)-1}((p,1_{\widehat{P}});\mathbb{K})\neq 0\}=\{p_1, p_2\}.
\end{align*}
Summing up, we have that our bounds in this case are
\[
-\infty=\reg(K^0 (R/I))<0,\ 0=\reg (K^1 (R/I))\leq 0<1,\ 2=\reg (K^2 (R/I))\leq 2.
\]
This shows, in particular, that our bound for $\reg (K^1 (R/I))$ is better than Kumini--Murai's one, even in the monomial case.
\end{ex}
Our next example in the monomial case illustrates that in our upper bound for the regularity it is not only important the poset of sums, but also the dimensions of the ideals on it.

\begin{ex}\thlabel{non equdimensional example}
Let $I=(xy,xz)\subset R:=\mathbb{K}[x,y,z].$ In this case, we know that $I=(x)\cap (y,z),$ and that our poset has vertices
\[
(x)=p_1,\ (y,z)=p_2,\ (x,y,z)=p_3.
\]
Moreover, the deficiency modules $K^j (R/I)$ are non--zero for $j=1$ and $j=2.$ Indeed, $R/I$ is not Cohen--Macaulay because it is actually non--pure. Actually, in this case, whereas the simplicial complex attached to $I$ through the Stanley--Reisner correspondence is not Cohen--Macaulay, the order complex attached to our poset $P$ is Cohen--Macaulay. Now, it only misses to compute the sets $S_j$.

First of all, notice that
\begin{align*}
& \{p\in P:\ 0\geq\dim (R/I_p)\}=\{p_3\},\ \{p\in P:\ 1\geq\dim (A/I_p)\}=\{p_2,\ p_3\},\\
& \{p\in P:\ 2\geq\dim (A/I_p)\}=P.
\end{align*}
Moreover, we also have $(p_1,1_{\widehat{P}})=\emptyset=(p_2,1_{\widehat{P}}),\ (p_3,1_{\widehat{P}})=\{p_1\}\cup\{p_2\}$. This implies that
\begin{align*}
& \widetilde{H}_{j-\dim (R/I_{p_1})-1}((p_1,1_{\widehat{P}});\mathbb{K})\neq 0\Longleftrightarrow j=2,\\
& \widetilde{H}_{j-\dim (R/I_{p_2})-1}((p_2,1_{\widehat{P}});\mathbb{K})\neq 0\Longleftrightarrow j=2,\\
& \widetilde{H}_{j-\dim (R/I_{p_3})-1}((p_3,1_{\widehat{P}});\mathbb{K})\neq 0\Longleftrightarrow j=1.
\end{align*}
This shows that
\begin{align*}
& S_0=\{p\in P:\ 0\geq\dim (A/I_p),\ \widetilde{H}_{0-\dim(A/I_p)-1}((p,1_{\widehat{P}});\mathbb{K})\neq 0\}=\emptyset,\\
& S_1=\{p\in P:\ 1\geq\dim (A/I_p),\ \widetilde{H}_{1-\dim(A/I_p)-1}((p,1_{\widehat{P}});\mathbb{K})\neq 0\}=\{p_2, p_3\},\\
& S_2=\{p\in P:\ 2\geq\dim (A/I_p),\ \widetilde{H}_{2-\dim(A/I_p)-1}((p,1_{\widehat{P}});\mathbb{K})\neq 0\}=\{p_1\}.
\end{align*}
Summing up, we have that our bounds in this case are
\[
-\infty=\reg(K^0 (R/I))<0,\ 0=\reg (K^1 (R/I))=0,\ 2=\reg (K^2 (R/I))=2.
\]
\end{ex}

Now, we want to illustrate Theorem \ref{deficiency in case of toric face rings} by means of the following example, which is given by a non Cohen--Macaulay toric face ring borrowed from \cite[Example 4.2.3]{HopNguyenthesis}.

\begin{ex}\thlabel{a non C.M. toric face ring}
We consider the toric face ring
\[
R:=\frac{\mathbb{Q}[x_1,x_2,x_3,x_4,x_5,x_6]}{(x_1x_2-x_4^2,x_2x_3-x_5^2,x_4x_6,x_5x_6)}.
\]
In this case, this is a three dimensional ring and its unique non--zero deficiency modules are $K^2 (R)$ and $K^3 (R)$ with respective regularities $0$ and $3.$ Moreover, we have that
\[
I:=(x_1x_2-x_4^2,x_2x_3-x_5^2,x_4x_6,x_5x_6)=(x_2x_3-x_5^2,x_1x_2-x_4^2,x_6)\cap (x_1,x_3,x_4,x_5)\cap (x_2,x_4,x_5).
\]
is the prime decomposition of the defining ideal of $R.$ In this way, the poset $P$ of all the sums of the prime components of $I$ turns out to be
\begin{align*}
& p_1=(x_2x_3-x_5^2,x_1x_2-x_4^2,x_6),\ p_2=(x_1,x_3,x_4,x_5),\ p_3=(x_2,x_4,x_5),\\
& p_4=(x_1,x_3,x_4,x_5,x_6),\ p_5=(x_2,x_4,x_5,x_6),\ p_6=(x_1,x_2,x_3,x_4,x_5),\\
& p_7=(x_1,x_2,x_3,x_4,x_5,x_6),
\end{align*}
In this case, since
\begin{align*}
& S_2=\{p\in P:\ 2\geq\dim (A/I_p),\ \widetilde{H}_{2-\dim(A/I_p)-1}((p,1_{\widehat{P}});\mathbb{K})\neq 0\}=\{p_2,p_4,p_6,p_7\},\\
& S_3=\{p\in P:\ 3\geq\dim (A/I_p),\ \widetilde{H}_{3-\dim(A/I_p)-1}((p,1_{\widehat{P}});\mathbb{K})\neq 0\}=\{p_1,p_3,p_5\},
\end{align*}
we can guarantee, summing up, that our bound in this case is
\[
0=\reg (K^2 (R))\leq 2=2,\ 3=\reg (K^3(R))\leq 3.
\]
This example shows that in general the inequality
\[
\reg(K^j(A/I))\leq\max_{p\in S_j} (\dim (A/I_p))
\]
might be strict even in case $K^j (A/I)\neq 0$.
\end{ex}

We conclude with an example of a binomial edge ideal borrowed from \cite[Example 3.2]{AlvarezMontaner2020}.

\begin{ex}
Let $A:=\mathbb{Q}[x_1,\ldots,x_5,y_1,\ldots,y_5],$ and let $J_G$ be the binomial edge ideal attached to the $5$--th path
\[
G:=\{\{1,2\},\{2,3\},\{3,4\},\{4,5\}\}.
\]
In this case, this is a six dimensional ring and its unique non--zero deficiency module $K^6 (R)$ has regularity $6.$ In this case, the poset $Q_{J_G}$ has $17$ elements and the regularities of the corresponding canonical modules of this list are (without repetitions)
\[
\{3,4,5,6\}.
\]
%In this way, the upper bound obtained in Theorem \ref{deficiency in case of binomial edge ideals} boils down in this case to an equality.
\end{ex}

Finally, our goal is to illustrate our results for binomial edge ideals with the below example, which is a non Cohen--Macaulay ring attached to a binomial edge ideal given by a complete bipartite graph. For more information about this specific kind of binomial edge ideals attached to complete bipartite graphs the reader can consult the work by Schenzel and Zafar, specially \cite[Section 4]{SchenzelZafarbinomialedgeideals}.

\begin{ex}
Let $A:=\mathbb{Q}[x_1,\ldots,x_8,y_1,\ldots,y_8],$ and let $J_G$ be the binomial edge ideal attached to the complete graph $K_{2,3}$. In this case, this is a ten dimensional ring and its unique non--zero deficiency modules are $K^j (R),$ where $j\in\{5,6,7,9,10\},$ with respective regularities $\{4,6,6,9,10\}$. In this case, the poset $Q_{J_G}$ has $6$ elements $\{p_1,\ldots, p_6\}$ such that
\[
\dim (A/I_{p_1})=10,\ \dim (A/I_{p_2})=9,\ \dim (A/I_{p_3})=6,\ \dim (A/I_{p_4})=6,\ \dim (A/I_{p_5})=0,\ \dim (A/I_{p_6})=4.
\]
The poset $Q_{J_G}$ in this specific example can be represented through the following Hasse--Voght diagram.
\[
\xymatrix{I_{p_1}&\hbox{}&  I_{p_2}& \hbox{}& I_{p_4}\\ \hbox{}& I_{p_3}\ar@{-}[ul]\ar@{-}[ur]& \hbox{}& I_{p_6}\ar@{-}[ul]\ar@{-}[ur]& \hbox{}\\ \hbox{}& \hbox{}&  I_{p_5}\ar@{-}[ul]\ar@{-}[ur]& \hbox{}& \hbox{}}
\]
Now, we want to write down the sets $S_j$'s for the non--zero canonical modules of $A/J_G$. We have
\begin{align*}
& S_5=\{p\in P:\ 5\geq\dim (A/I_p),\ \widetilde{H}_{5-\dim(A/I_p)-1}((p,1_{\widehat{P}});\mathbb{K})\neq 0\}=\{p_6\},\\
& S_6=\{p\in P:\ 6\geq\dim (A/I_p),\ \widetilde{H}_{6-\dim(A/I_p)-1}((p,1_{\widehat{P}});\mathbb{K})\neq 0\}=\{p_4\},\\
& S_7=\{p\in P:\ 7\geq\dim (A/I_p),\ \widetilde{H}_{7-\dim(A/I_p)-1}((p,1_{\widehat{P}});\mathbb{K})\neq 0\}=\{p_3\},\\
& S_9=\{p\in P:\ 9\geq\dim (A/I_p),\ \widetilde{H}_{9-\dim(A/I_p)-1}((p,1_{\widehat{P}});\mathbb{K})\neq 0\}=\{p_2\},\\
& S_{10}=\{p\in P:\ 10\geq\dim (A/I_p),\ \widetilde{H}_{10-\dim(A/I_p)-1}((p,1_{\widehat{P}});\mathbb{K})\neq 0\}=\{p_1\}.\\
\end{align*}
On the one hand, since $\dim (A/I_{p_6})=4$, we can guarantee that $4=\reg (K^5 (A/J_G))\leq 4<5$. On the other hand, since $\dim (A/I_{p_3})=6$, we can ensure that $6=\reg (K^7(A/J_G))\leq 6<7$. Hence once again our bound for this example turns out to be better than $\reg (K^j (A/J_G))\leq j$.
\end{ex}

\section*{Acknowledgements}
The authors would like to thank Josep \`Alvarez Montaner and Fernando Sancho de Salas for several interesting comments concerning the contents of this paper. The authors received partial support by grant PID2022-137283NB-C22 funded by  MICIU/AEI/10.13039/501100011033.


\bibliographystyle{alpha}
\bibliography{AFBoixReferences}

\end{document}

\item Finally, in case $A/I$ is sequentially Cohen--Macaulay, we can ensure that $S_j\neq\emptyset$ if and only if $\depth (A/I)\leq j\leq\dim (A/I)$.

\begin{disc}
Keeping the notations and assumptions of the above Example, in fact the spectral sequence obtained in \thref{sucesion espectral y colapso todo en uno: segunda parte} boils down, in case $T=\hbox{}^*\Hom_A (-,A(-n))$ to the below spectral sequence:
\[
E_2^{-i,j}=\bigoplus_{j=h_q} \hbox{}^*\Ext_A^{h_q} \left(V_q,A(-n)\right)^{\oplus m_{i,q}}\xymatrix{ \ar@{=>}[r]_-{i}& }\hbox{}^*\Ext_A^{j-i} \left(\lim_{p\in P}V_p,A(-n)\right),
\]
where $m_{i,q}:=\dim_{\K} \left(\widetilde{H}_{i-1} \left(\left(q,1_{\widehat{P}}\right);\K\right)\right)$.
\end{disc}

\begin{rk}
As consequence of \cite[Proposition 4.1]{DaoMaVarbaro}, we know that if $\mathbb{K}$ is a field of prime characteristic such that $\mathbb{K}[\mathcal{M}_{\Sigma}]$ is $F$--pure (see \cite[Theorem 6.3]{HopNguyen2012}) then we have that
\[
\reg(K^j(B))\leq\dim (K^j (B))\leq j.
\]
This shows, in particular, that the upper bound we obtain in \thref{deficiency in case of toric face rings} also holds for other classes of toric face rings beyond the case we cover.
\end{rk}

\begin{ex}\thlabel{an example of a genuine toric face ring}
We consider the toric face ring
\[
R:=\frac{\mathbb{Q}[x,y,z,w]}{(xy-w^2,zw)}.
\]
In this case, this is a two dimensional ring and its unique non--zero deficiency module $K^2 (R)$ has regularity $2.$ Moreover, we have that
\[
I:=(xy-w^2,zw)=(x,w)\cap (y,w)\cap (xy-w^2,z)
\]
is the prime decomposition of the defining ideal of $R.$ In this way, the poset $P$ of all the sums of the prime components of $I$ turns out to be
\[
\{(x,w),(y,w),(xy-w^2,z),(x,y,w),(x,z,w),(y,z,w),(x,y,z,w)\},
\]
and the regularities of the corresponding canonical modules of this list are
\[
\{2, 2, 1, 2, 1, 1, 0\}.
\]
In this way, the upper bound obtained in Theorem \ref{deficiency in case of toric face rings} boils down in this case to an equality.
\end{ex}

We end up this part with a similar example of non--Cohen Macaulay toric face ring borrowed from \cite[Example 4.5]{Yanagawaseminormaltoricfacerings}.

\begin{ex}
We consider the toric face ring
\[
R:=\frac{\mathbb{Q}[x,y,z,u,v,w]}{(xv-uy,vz-ywxz-uw,uvw,uvz)}.
\]
In this case, this is a three dimensional ring and its unique non--zero deficiency modules are $K^2 (R)$ and $K^3 (R)$ with respective regularities $0$ and $3.$ Moreover, we have that
\[
I:=(xv-uy,vz-ywxz-uw,uvw,uvz)=(xv-yu,z,w)\cap (zv-yw,x,u)\cap (xz-uw,y,v).
\]
is the prime decomposition of the defining ideal of $R.$ In this way, the poset $P$ of all the sums of the prime components of $I$ turns out to be
\begin{align*}
& p_1=(w, z, yu - xv),\ p_2=(u, x, zv - yw),\ p_3=(w, u, z, x),\ p_4=(v, y, xz - uw),\\
& p_5=(w, v, z, y),\ p_6=(v, u, y, x),\ p_7=(x,y,z,u,v,w),
\end{align*}
and the regularities of the corresponding canonical modules of this list are
\[
\{3,3,2,3,2,2,0\}.
\]
\end{ex}

\section{The case of monomial ideals}\label{section: the case of monomial ideals}
The purpose of this section is to obtain upper bounds for the deficiency modules of monomial algebras. We obtain such upper bounds for monomial ideals inside a standardly graded polynomial ring. Indeed, our main result is the following statement, which is an immediate consequence of \thref{the expected bound when P is a C.M. poset}.

\begin{teo}\thlabel{deficiency in case of monomial ideals}
Let $\mathbb{K}$ be a field, let $A=\mathbb{K}[x_1,\ldots,x_n]$ be the polynomial ring in $n$ variables over the field $\mathbb{K},$ and let $I\subseteq R$ be a monomial ideal. Then, we have that $\reg(K^j(A/I))\leq j$.
\end{teo}

\section{Results involving other ultrametric functions}\label{results concerning other ultrametric functions}

The goal of this section is to specialize some of our results to other ultrametric functions different from the regularity. For instance, the following result provides an interesting upper bound for the depth of deficiency modules of rings of sections.

\begin{teo}\thlabel{main result for the depth}
Under the assumptions and notations of Discussion \ref{the inequality for any ultrametric function}, assume that $\mu=-\depth$ and that, for any $p\in P,$ $A/I_p$ is Gorenstein. Then, for any integer $j\geq 0$ we have
\[
\depth\left(K^j \left(\lim_{p\in P}A/I_p\right)\right)\geq j.
\]
\end{teo}

\begin{proof}
Fix $p\in P,$ since $A/I_p$ is Gorenstein, it is in particular Cohen--Macaulay and therefore we have that $K^{\dim (A/I_p)}(A/I_p)$ is a Cohen--Macaulay module, hence we have
\[
\depth (K^{\dim (A/I_p)}(A/I_p))=\dim (K^{\dim (A/I_p)}(A/I_p)).
\]
Moreover, since $A/I_p$ is Gorenstein we have that $K^{\dim (A/I_p)}(A/I_p)$ is abstractly isomorphic to $A/I_p,$ which implies that
\[
\dim (K^{\dim (A/I_p)}(A/I_p))=\dim (A/I_p).
\]
In this way, the result follows directly from Discussion \ref{the inequality for any ultrametric function} choosing $\mu$ as $-\depth,$ and keeping in mind that
\[
S_j=\{p\in P:\ j\leq\dim (A/I_p),\ m_{n-j,p}\neq 0\}\qedhere.
\]
\end{proof}

Our next result follows immediately by combining \thref{main result for the depth} with the Auslander--Buchsbaum formula. 

\begin{teo}\thlabel{main result for the pd}
Under the assumptions and notations of Discussion \ref{the inequality for any ultrametric function}, assume that $\mu=-\depth$ and that, for any $p\in P,$ $A/I_p$ is Gorenstein. Then, for any integer $j\geq 0$ we have
\[
\projdim\left(K^j \left(\lim_{p\in P}A/I_p\right)\right)\leq \dim (A)-j,
\]
where $\projdim$ refers to the projective dimension in the category of $A$--modules.
\end{teo}

The reader will easily note that our bound is not exactly the one obtained in \cite{KumminiMurai2011}, however we want to single out it because our approach and proof is completely independent from the one used in \cite{KumminiMurai2011}.

However, we can check (using Macaulay2) that the corresponding order complex attached to $P$ is Cohen--Macaulay.

Finally, performing similar calculations we already did in previous examples we can check that
\[
S_2=\{p_1,p_2,p_3,p_4,p_5,p_6\},\ S_3=\{p_1,p_2,p_4\},
\]
and therefore in this case we also have that $\reg (K^j(R))\leq j.$

However, we can check (using Macaulay2) that the corresponding order complex attached to $P$ is Cohen--Macaulay. Finally, performing similar calculations we already did in previous examples we can check that
\[
S_2=\{p_1,p_2,p_3\},\ S_3=\{p_1,p_3\},
\]
and therefore in this case we also have that $\reg (K^j(R))\leq j.$

In this case. the Hasse--Voght diagram of our poset can be depicted as follows:
\[
\xymatrix{\hbox{}& I\\ I_{p_1}\ar@{.}[ur]& I_{p_2}\ar@{.}[u]& I_{p_3}\ar@{.}[ul]\\ I_{q_1}\ar@{-}@{-}[u]\ar@{-}[ur]& \hbox{}& I_{q_2}\ar@{-}[ul]\ar@{-}[u]\\ & I_{q_3}\ar@{-}[ul]\ar@{-}[ur]& }
\]
We can check that the corresponding order complex $\Delta (Q_{J_G})$ is Cohen--Macaulay.

Before giving more concrete examples, we recall the following notions.

\begin{df}\thlabel{rank of a poset}
Let $Q$ be a poset. A \textbf{chain} of $Q$ is a totally ordered subset of $Q;$ moreover, the \textbf{rank of a chain} is the number of elements in it. On the other hand, we define the \textbf{rank} of $Q$ (denoted $\rank (Q)$) as the rank of the longest chain in $Q.$ Finally, we define the \textbf{length} of $Q$ as $\lambda (Q):=\rank(P)-1.$
\end{df}

\begin{df}\thlabel{Cohen--Macaulay posets}
Let $Q$ be an arbitrary poset, and let $\mathbb{K}$ be a field. We say that $Q$ is \textbf{Cohen--Macaulay with respect to $\mathbb{K}$} if its corresponding order complex $\Delta (Q)$ is Cohen--Macaulay with respect to $\mathbb{K}$ as a simplicial complex.
\end{df}

\subsection{Examples in the monomial case}

Next example shows, on the one hand, that the assumptions imposed in Corollary \ref{in the Sj's only maximal elements} are satisfied in at least one case; on the other hand, it also shows that the equidimensional assumption can not be removed in Proposition \ref{this is what happens when Sj is empty}.


Next example shows that, in general, even in the Stanley--Reisner case, we can not expect a relation between the Cohen--Macaulayness of the simplicial complex given by the Stanley--Reisner correspondence, and the order complex attached to our poset $P.$

\begin{ex}\thlabel{order complex not CM but Delta CM}
Let $I=(xy,xz,yz)\subset R:=\mathbb{K}[x,y,z].$ In this case, we know that $I=(x,y)\cap (x,z)\cap (y,z),$ and that our poset has vertices
\[
(x,y)=p_1,\ (x,z)=p_2,\ (y,z)=p_3,\ p_4=(x,y,z).
\]
Moreover, the deficiency modules $K^j (R/I)$ are non--zero for $j=1.$ Indeed, $R/I$ is Cohen--Macaulay in this case; on the other hand, we check through Macaulay2 that $P$ is not Cohen--Macaulay.
\end{ex}
Next example illustrate that the sets $S_j$ might be empty, as we have already shown through Proposition \ref{this is what happens when Sj is empty}.


However, checking whether $P$ is Cohen--Macaulay or not might be highly non--trivial, as the following example, borrowed from \cite{HooperVanTuylvdWcomplex}, shows.

\begin{ex}\thlabel{the van der Waerden complex}
Let $I\subset R:=\mathbb{Q}[x,y,z,w,a,b,c]$ be the squarefree monomial ideal attached to the so--called van der Waerden complex $\operatorname{vdW}(7,2);$ explicitly, we have
\[
I=(xyw, xzw, ya, xwa, xb, zb, yc, xzc, zwc, xac, wac, wbc).
\]
As proved in \cite{HooperVanTuylvdWcomplex}, this simplicial complex is pure but not Cohen--Macaulay in this case; indeed, the deficiency modules $K^j (R/I)$ are non--zero for $j=2$ and $j=3.$ On the other hand, our poset $P$ has $26$ vertices, so it does not seem easy neither to draw it, nor to check whether it is Cohen--Macaulay. Finally, we also want to mention that the Alexander dual $I^{\vee}$ is a squarefree monomial ideal such that $R/I^{\vee}$ is not Cohen--Macaulay, and the corresponding poset $P$ attached to $I^{\vee}$ has $67$ vertices.
\end{ex}

We include one more example where all the involved simplicial structures are not Cohen--Macaulay.

\begin{ex}\thlabel{nothing is Cohen--Macaulay here}
Let $I:=(xw, xz,xy, yzw)\subset\mathbb{Q}[x,y,z,w]:=R.$ First of all, in this case $I=I^{\vee};$ that is, $I$ agree with its Alexander dual. Moreover, $R/I$ is not Cohen--Macaulay, because $R/I$ is a two dimensional ring with $K^j (R/I)\neq 0$ for $j=1$ and $j=2.$ On the other hand, in this case our poset $P$ has the following eight vertices:
\begin{align*}
& p_1=(x,y),\ p_2=(x,z),\ p_3=(x,w),\ p_4=(y,z,w),\ p_5=(x,y,z),\\
& p_6=(x,y,w),\ p_7=(x,z,w),\ p_8=(x,y,z,w).
\end{align*}
We have checked using Macaulay2 that the corresponding order complex $\Delta (P)$ attached to $P$ is not Cohen--Macaulay. Moreover, looking at the sets $S_1$ and $S_2$ we have that
\[
S_1=\{p_4\},\ S_2=\{p_1,p_2,p_3\}
\]
Indeed, let us justify the calculation of $S_1;$ first of all, we observe that
\[
1\leq\dim (A/I_p)\Longleftrightarrow p\neq p_8.
\]
On the other hand, we need to distinguish two cases; on the one hand, if $\dim (A/I_p)=2$ then $p\in\{p_1,p_2,p_3\},$ hence $(p,1_{\widehat{P}})=\emptyset$ and therefore
\[
\widetilde{H}_{2-1-1}((p,1_{\widehat{P}});\mathbb{K})=\widetilde{H}_0 (\emptyset;\mathbb{K})=0.
\]
On the other hand, if $\dim (A/I_p)=1$ then $p\in\{p_4,p_5,p_6,p_7\}.$ However, for $p\in\{p_5,p_6,p_7\}$ we have that $(p,1_{\widehat{P}})\neq\emptyset$ and therefore, in this case,
\[
\widetilde{H}_{1-1-1}((p,1_{\widehat{P}});\mathbb{K})=\widetilde{H}_{-1}((p,1_{\widehat{P}});\mathbb{K})=0.
\]
Joining these two facts we finally conclude that $S_1=\{p_4\},$ as claimed. We also want to mention that our calculation also implies that $\reg (K^1 (R/I))=0<1$ and $\reg (K^2 (R/I))=2.$
\end{ex}

\begin{rk}
It is known, by \cite[Theorem 1.1]{KumminiMurai2011} that, under the assumptions of \thref{deficiency in case of monomial ideals}, we have that
\[
\reg(K^j(A/I))\leq\dim (K^j (A/I)),
\]
and since $\dim (K^j (A/I))\leq j$ \cite[Lemma 1.9 (a)]{SchenzelUAB}, we finally obtain
\[
\reg(K^j(A/I))\leq j,
\]
which is exactly \cite[Corollary 1.2]{KumminiMurai2011}. The new aspect of the upper bound got in Theorem \ref{deficiency in case of monomial ideals} is not its statement but its proof, because it is clear that, in general, Kumini--Murai upper bound is sharper than the one we obtain here (see however \thref{sometimes S_j is empty}). Remember that the inequality
\[
\reg(K^j(A/I))\leq\dim (K^j (A/I)),
\]
is also true if $X=\operatorname{Proj}(A/I)$ defines either a Du Bois singularity in characteristic zero \cite[Proposition 3.1]{DaoMaVarbaro} or an $F$--pure singularity in prime characteristic \cite[Proposition 4.1]{DaoMaVarbaro}, so similar comments apply in these settings.
\end{rk}

As immediate consequence of \thref{the expected bound when P is a C.M. poset} we obtain the following statement.

\begin{cor}\thlabel{in the Sj's only maximal elements}
Preserving the notations and assumptions of \thref{the expected bound when P is a C.M. poset}, let $j\geq 0$ be such that $K^j (R/I)\neq 0,$ and assume that $S_j\subseteq\{p\in P:\ j=\dim (A/I_p)\}.$ Then, $\reg (K^j (R/I))\leq j.$
\end{cor}
Using our formalism, we can prove the following result.

\begin{prop}\thlabel{this is what happens when Sj is empty}
Let $\mathbb{K}$ be a field, let $A=\mathbb{K}[x_1,\ldots,x_n]$ be the polynomial ring in $n$ variables over the field $\mathbb{K},$ and let $I\subseteq R$ be a homogeneous ideal. Assume that $I$ is equidimensional, and that the Lyubeznik graph (more precisely, the $1$--skeleton of the simplicial complex defined in \cite[page 622]{Lyu2006}) has no edges. Finally, we assume that $\depth (A/I)\geq 1.$ Then, we have, for each integer $j\geq 1$ with $K^j (A/I)\neq 0,$ that $\reg (K^j (A/I))\leq j.$
\end{prop}

\begin{proof}
Set $d:=\dim (A/I)\geq 1$ and fix $j\geq 1$ with $K^j (R/I)\neq 0.$ Since $I$ is equidimensional, any minimal prime of $I$ (that correpond in our poset $P$ to vertices $p_1,\ldots, p_r$) has dimension exactly $d;$ on the other hand, our assumption on the Lyubeznik graph implies that
\[
\dim (A/I_p)=0
\]
for each $p$ that do not correspond to a minimal prime of $I.$ This implies that
\[
S_j=\begin{cases}
\emptyset,\text{ if }j<d,\\
\{p_1,\ldots,p_r\},\text{ if }j=d.
\end{cases}
\]
This shows, in particular, that
\[
\reg (K^j (A/I))\leq\begin{cases}
0,\text{ if }j<d,\\
d,\text{ if }j=d,
\end{cases}
\]
which in particular proves our statement.
\end{proof}

\item If, in addition, $P$ is a Cohen--Macaulay poset with respect to $\mathbb{K},$ then we have that
\[
\reg(K^j(A/I))\leq j+\max_{p\in S_j}(\rank(p,1_{\widehat{P}})).
\]


Assume now that $P$ is Cohen--Macaulay, this is equivalent, by definition, to say that $\Delta (P)$ is Cohen--Macaulay as simplicial complex and therefore, for any $p\in P,$ we have, applying Reisner criterion \cite[Theorem 1]{Reisnercriterion} that
\[
\widetilde{H}_{t} ((p,1_{\widehat{P}});\K)=\widetilde{H}_{t} (\link_{\Delta(P)} (p);\K)=0
\]
for all $t\leq\dim (\link_{\Delta(P)} (p))-1=\lambda((p,1_{\widehat{P}}))-1.$ (Remember that $\link_{\Delta(P)} (p)$ is isomorphic to $(p,1_{\widehat{P}})$ viewed as poset, see either \cite[second paragraph of page 229]{BaclawskiCMposets} or \cite[Lemma 4.9]{SanchoDeSalasStanleyReisnertheory}). This implies that
\[
\widetilde{H}_{\dim (A/I_p)-j-1}((p,1_{\widehat{P}});\K)=0
\]
for all $\dim (A/I_p)-j-1\leq\lambda((p,1_{\widehat{P}}))-1.$ In this way, we have that
$\widetilde{H}_{\dim (A/I_p)-j-1}((p,1_{\widehat{P}});\K)\neq 0$ if and only if $\dim (A/I_p)=j+\dim (\link_{\Delta(P)} (p))+1=j+\lambda((p,1_{\widehat{P}}))+1=j+\rank(p,1_{\widehat{P}}),$ which in turn implies that $S_j=\{p\in P:\ \dim(A/I_p)=j+\rank(p,1_{\widehat{P}})\}.$ In this way, we finally obtain the upper bound
\[
\reg(K^j(A/I))\leq j+\max_{p\in S_j}(\rank(p,1_{\widehat{P}})),
\]
as claimed.

On the other hand, in Section \ref{results concerning other ultrametric functions} we apply some of our results to other ultrametric functions different from the regularity.

%\begin{teo}\label{formula de Hochster: no me digas: segunda parte}
%Under the assumptions of \textbf{Theorem} \ref{sucesion espectral y colapso
%todo en uno: segunda parte}, there is a $\K$-vector space
%isomorphism
%\[
%\R^j T\left(\lim_{p\in P} V_p\right)\cong\bigoplus_{q\in P} \R^{h_q}
%T (V_q)^{\oplus m_{j,q}}.
%\]
%where $m_{j,q}=\dim_{\K} \widetilde{H}_{j-h_q-1} ((q,1_{\widehat{P}});
%\K).$ More generally, if $\mathcal{B}$ is an abelian, semisimple category containing the category of $A$--modules such that $T(\mathcal{B})\subseteq\mathcal{B}$, and if $V\in\Inv(\widehat{P},\mathcal{B})$, then the previous bijection is actually an isomorphism in the category $\mathcal{B},$
%\end{teo}


\begin{lm}\label{contractibilidad de abiertos en la topologia de Alexandrov}
Let $(P,\leq)$ be a poset regarded as a topological space with the Alexandrov topology. Then, for any $p\in P$,
the basic open set $[p,1_{\widehat{P}})$ is contractible.
\end{lm}

Now, the case of projective objects; we refer to \cite[Satz 2 and Beweis]{Nobeling1962} for details.

\begin{teo}\label{proyectivos que interesan}
Given a projective $A$-module $F$ and given $p\in P$, the inverse system $F_{\leq p}$ defined by
\[
\left(F_{\leq p}\right)_q:=\begin{cases} F,\text{ if }q\in (0_{\widehat{P}},p],\\ 0,\text{otherwise}\end{cases}
\]
is a projective object of $\Inv (P,\cA)$; in particular, $A_{\leq p}$ is so. Here, given $q\leq r$ the homomorphism $\xymatrix@1{\left(F_{\leq p}\right)_r\ar[r]& \left(F_{\leq p}\right)_q}$
is zero if either source or target is zero, otherwise it equals the identity on
$F.$ Furthermore, any inverse system $G$ admits a projective (not necessarily minimal) resolution
\[
\xymatrix{\ldots\ar[r] & F_1\ar[r]& F_0\ar[r]& G\ar[r]& 0}
\]
such that, for any $i\in\N$, $F_i$ can be expressed as a direct sum of projectives of the form $A_{\leq p}$.
\end{teo}

The following lemma will provide the abutment of the spectral sequence we want to construct.

\begin{lm}\label{lema clave en este arreglo}
Let $G$ be a projective object of $\Inv (P,\cA)$ of the form
\[
G=\bigoplus_{j\in J} A_{\leq q_j},
\]
where $J$ is a finite index set and $q_j\in P$ (see \textbf{Theorem} \ref{proyectivos que interesan}). Then, the augmented
chain complex
\[
\roos_* \left(\mathcal{T}\left(G\right)\right)\longrightarrow\left(\colim_{p\in P}\circ\mathcal{T}\right)
\left(G\right)\longrightarrow 0
\]
is exact.
\end{lm}

\begin{proof}
Since $\roos_*$, $\mathcal{T}$ and the colimit functor commutes with finite direct sums we may suppose, without
loss of generality, that $G=A_{\leq q}$ for some fixed $q\in P$; regardless, in this case, our augmented chain
complex is exactly the one for computing the simplicial homology of the interval $[q,1_{\widehat{P}})$ (indeed,
notice that we have to take this interval because $\mathcal{T}$ is contravariant) with coefficients in $Z$. But
$[q,1_{\widehat{P}})$ is contractible by \textbf{Lemma} \ref{contractibilidad de abiertos en la topologia de Alexandrov};
the proof is therefore completed.
\end{proof}

\begin{lm}\label{playing with associated primes}
Let $(A,\mathfrak{m})$ be a commutative Noetherian local ring where there exists a canonical
module $\omega_A,$ (so, $A$ is a Cohen--Macaulay ring
which can be expressed as homomorphic image of a Gorenstein
local ring $B$ of dimension $d$) let $I_p,\ I_q$ be two
ideals of $A$ of finite projective dimension
such that $I_q\not\subseteq\mathfrak{p}$ for any
prime ideal $\mathfrak{p}\supseteq I_p$ such that
$\dim (A/\mathfrak{p})=d_p.$ Then,
$\Hom_A (\Ext_A^{d-d_q} (A/I_q,\omega_A),
\Ext_A^{d-d_p} (A/I_p,\omega_A))=0.$
\end{lm}

We need a graded version of the above result; namely:

\begin{lm}
Let $\mathbb{K}$ be a field, let $A=\mathbb{K}[x_1,\ldots,x_n]$ be a standardly graded polynomial ring over $\mathbb{K},$ and let $I_p,\ I_q$ be two ideals of $A$ such that $I_q\not\subseteq\mathfrak{p}$ for any
prime ideal $\mathfrak{p}\supseteq I_p$ such that $\dim (A/\mathfrak{p})=d_p.$ Then,
$\hbox{}^*\Hom_A (\hbox{}^*\Ext_A^{n-d_q} (A/I_q,A(-n)),
\hbox{}^*\Ext_A^{n-d_p} (A/I_p,A(-n)))=0.$
\end{lm}

\subsection{Monomial ideals inside a multigraded polynomial ring}
Throughout this subsection (and only here), let $G$ be a finitely generated, free abelian group, let $\mathbb{K}$ be a field, let $A:=\mathbb{K}[x_1,\ldots,x_n]$ be a polynomial ring over $\mathbb{K}$ graded by $G,$ and let $R:=A/I,$ where $I$ is a monomial ideal. Given a finitely generated, graded $A$--module $M$ we denote by $\reg(M)$ the multigraded regularity as defined in \cite[Definition 4.1]{MaclaganSmith2004}. Our main result here is given by the following (completely general) upper bound for the multigraded regularity of the deficiency modules of $A$ modulo a monomial ideal.

\begin{teo}\label{multigraded regularity of deficiency modules}
For each $1\leq i\leq n,$ set $a_i:=\deg (x_i)$ and
\[
a:=\sum_{i=1}^n a_i.
\]
If $K^j(A/I):=\hbox{}^* \Ext_A^{n-j} (A/I,A(-a)),$ then we have that
\[
\bigcap_{p\in S_j}(\reg(K^j (A/I_p)))\subseteq\reg(K^j(A/I)),
\]
where
\[
S_j:=\{p\in P:\ j\leq \dim (A/I_p),\ m_{n-j,p}\neq 0\}.
\]
\end{teo}

\section{The case of rings realized by retracts}

The goal of this section is to show that, for the case of a graded ring realized by retracts, it is easier to obtain our bounds for the regularity. Let us recall first this notion, following \cite[Definition 1.1]{Berkeschringsrealizedbyretracts}.

\begin{df}\label{rings realized by retracts}
Let $\mathbb{K}$ be a field, let $R$ be graded $\mathbb{K}$--algebra, and let $\{S_l\}_{l\in\Lambda}$ be a finite collection of domains that are algebra retracts of $R$ with defining graded $\mathbb{K}$--algebra homomorphisms $i_l:\ \xymatrix@1{S_l\ar@{^(->}[r]& R}$ and $\pi_l:\ \xymatrix@1{R\ar@{->>}[r]& S_l}$ such that $\pi_l\circ i_l=1_{S_l}.$ For each $l,$ set $\mathfrak{p}_l:=\ker (\pi_l).$ We call $R$ a $\mathbb{K}$--algebra \textbf{realized by retracts} $\{S_l\}_{l\in\Lambda}$ if the following two conditions are satisfied.

\begin{enumerate}[(i)]

\item The natural map
\[
\xymatrix{R\ar[r]^-{\phi}& }\bigoplus_{l\in\Lambda} S_l
\]
given by $r\longmapsto (\pi_l (r))$ is injective.

\item For each $i\in\Lambda,$ we have
\[
\bigcap_{\substack{l\in\Lambda\\ i\neq l}}\mathfrak{p}_l\neq (0).
\]

\end{enumerate}
\end{df}
For graded algebras realized by retracts, the following result is inmmediate.

\begin{prop}
Let $\mathbb{K}$ be a field,
\[
R=\bigoplus_{n\geq 0} R_n
\]
be a positively graded ring such that $R=\mathbb{K} [R_1],$ and let $\mathcal{A}$ be the category of finitely generated, graded $R$--modules. Let $\mu$ be an ultrametric function on $\mathcal{A},$ and assume that $R$ is realized by retracts $\{S_l\}_{l\in\Lambda}.$ Then, we have that
\[
\mu (R)\leq\max_{l\in\Lambda} \mu (S_l).
\]
\end{prop}


On the other hand, in this case the poset $Q_{J_G}$ has six vertices that, following the notations given in \cite[Example 3.10]{AlvarezMontaner2020}, we denote by $p_1,\ p_2,\ p_3,\ q_1,\ q_2,\ q_3.$ We have that
\[
S_5=\{q_2\},\ S_6=\{p_1\},\ S_7=\{q_2\},\ S_9=\{p_2\},\ S_{10}=\{p_3\}.
\]
In this way, what we really obtain are the (in)equalities
\[
4<6,\ 6=6,\ 
\]

Actually, the upper bound obtained in Theorem \ref{deficiency in case of monomial ideals} is exactly the one obtained in \cite[Corollary 1.2]{KumminiMurai2011} if we put an additional assumption; namely:

\begin{teo}\label{deficiency in case of monomial ideals with extra assumption}
Let $\mathbb{K}$ be a field, let $A=\mathbb{K}[x_1,\ldots,x_n]$ be the polynomial ring in $n$ variables over the field $\mathbb{K},$ and let $I\subseteq R$ be a monomial ideal. Moreover, let $\Delta$ be te simplicial complex attached to the squarefree monomial ideal $\sqrt{I}.$ If $\sqrt{I}$ admits a linear resolution, then we have that
\[
\reg(K^j(A/I))\leq\max_{p\in S_j}(\dim (A/I_p))\leq j,
\]
where
\[
S_j:=\bigcup_{k\geq 0}S_{j,k},
\]
and
\[
S_{j,k}:=\{p\in P:\ j+k=\dim (A/I_p),\ m_{n-j,p}\neq 0\}.
\]
\end{teo}

\begin{proof}
The proof is the same; the fact that we obtain in this case a better upper bound is as follows. Let $p\in P,$ which corresponds to a non--face $F$ of $\Delta.$ Since $\sqrt{I}$ admits a linear resolution, we have by \cite{EagonReinerAlexanderduality} that $\Delta^{\vee}$ is Cohen--Macaulay and therefore we have that
\[
\widetilde{H}^t (\link_{\Delta^{\vee}}(F);\mathbb{K})=0\text{ for all }t<\dim (\link_{\Delta^{\vee}}(F)).
\]
In this way, this fact implies that
\[
S_j=\{p\in P:\ j=\dim (A/I_p)\}
\]
and therefore, in particular,
\[
\max_{p\in S_j}(\dim (A/I_p))=j,
\]
as claimed.
\end{proof}


\section{Sequential modules}
The goal of this section is to introduce the concept of sequentially $P$--modules that recover and generalize, for instance, sequentially Cohen--Macaulay modules.

\begin{df}\label{sequentially P modules}
Let $\mathbb{K}$ be a field, let $R=\mathbb{K}[x_1,\ldots ,x_n],$ standardly graded so that $\deg (x_i)=1$ for any $1\leq i\leq n,$ and let $M$ be a finitely generated, graded $R$--module. Given an integer $i\geq 0$ and given $\mathcal{C}_i$ some subcategory of the category of finitely generated, graded $R$--modules, we say that $M$ is \textbf{sequentially $P$} if there exists a filtration
\[
F_{\bullet}:\ 0=F_{-1}\subseteq F_0\subseteq F_1\subseteq\ldots\subseteq F_t=M
\]
such that, for each $0\leq i\leq t,$ $F_i/F_{i-1}$ is either zero or belongs to $\mathcal{C}_i$
\end{df}

\begin{ex}
At this point, some examples of the above notion are written down for the convenience of the reader.

\begin{enumerate}

\item If $\mathcal{C}_i$ is the category of Cohen--Macaulay modules of dimension $i,$ then, choosing $F$ as the dimension filtration of $M,$ we recover Stanley's notion of sequentially Cohen--Macaulay modules.

\item Given an integer $k\geq 0$, if $\mathcal{C}_i$ is the category of $k$--Buchsbaum modules of dimension $i,$ then, choosing $F$ again as the dimension filtration, we recover sequentially $k$--Buchsbaum modules as introduced in \cite[Definition 1]{HaHoa2008}.

\item If $\mathcal{C}_i$ is the category of $i$--reflexive modules, then we recover sequentially reflexive modules as introduced in \cite{Tonolo2004}.

\end{enumerate}
\end{ex}

\section{Unplugged: desired results}

\begin{enumerate}[(i)]

\item How to use \cite[Proposition 4.6]{AlvarezMontaner2020}.

\item The methods used here to obtain the regularity of some deficiency modules are different and independent from the ones used in \cite[Corollary 7]{BrodmannLinhSeiler2013}.

\item The methods used here to obtain the regularity of some deficiency modules are different and independent from the ones used in \cite[Theorem 9 and Theorem 14]{HoaHyry2006}.

\item The methods used here are different and independent from the ones used in \cite[Theorem 2.4 and Theorem 3.5]{ChardinHaHoa2011} to obtain the regularity of some Ext modules.

\end{enumerate}

\item What about obtaining lower bounds of the so--called multigraded regularity set as defined in \cite{MaclaganSmith2004}? See also \cite[Theorem 4.1]{MaclaganSmith2005}.

\item Notice that, by \cite[Proposition 2.3]{HaHoa2008}, if $R=\mathbb{K}[x_1,\ldots ,x_n]/I$ is a Cohen--Macaulay, graded ring, then
\[
\reg (K_R)=\dim (R).
\]
In particular, this shows that, for any graded $R$--module of the form
\[
M=\bigoplus_{i=1}^t K_R (-b_i),
\]
one has that
\[
\reg (M)=\dim (R)+\max_{1\leq i\leq t} (b_i).
\]
%Can we use this equality to derive some upper bound for the regularity of at least some Gorenstein modules \cite{SharpGorensteinmodules}?

\item The methods used here to obtain the regularity of some deficiency modules are different and independent from the ones used in \cite{HaHoa2008}.

\item In case $R$ is a polynomial ring, and $I$ is a monomial ideal, we have to recover \cite[Theorem 1.1]{KumminiMurai2011} or at least \cite[Corollary 1.2]{KumminiMurai2011} with an independent approach.

\item The same business works for a monomial ideal in a regular sequence.

in this case, the result
follows immediately combining (i) jointly with \cite[1.2.1 and 1.2.2]{HerzogHibi2011} and
our assumption that $S_j=\{b_j\}.$

Finally, set $\mathfrak{a}:=\sqrt{I},$ and assume that $J_e\neq 0.$ Up to relabeling we
can assume, without loss of generality, that there are $1\leq l<r<d$ integers such that
$\mathbf{x}^{\gamma}:=(x_1^{b_1}\cdots x_l^{b_l})^{p^e} (x_{l+1}^{b_{l+1}}\cdots x_r^{b_r})^{p^e-1}
\in J_e.$ Set
\[
m:=(x_1\cdots x_l)^{p^e} (x_{l+1}\cdots x_r)^{p^e-1}.
\]
We claim that $m\in\mathfrak{b}_e;$ first of all, assume that $m\in\mathfrak{a}^{[p^e]},$ so there
is a squarefree monomial generator $M\in\mathfrak{a}$ such that $M^{p^e}$ divides $m,$ hence we can
assume that there is $1\leq j_1<\ldots <j_s\leq l$ such that $M=x_{j_1}\cdots x_{j_s}.$ However, combining
our assumption on the sets $S_j$'s jointly with \cite[1.2.4]{HerzogHibi2011} one has that
$M=\sqrt{\widetilde{M}},$ where $\widetilde{M}:=x_{j_1}^{b_{j_1}}\cdots x_{j_s}^{b_{j_s}}\in I,$ and
therefore $\widetilde{M}^{p^e}$ divides $\mathbf{x}^{\gamma},$ a contradiction because
$\mathbf{x}^{\gamma}\notin I^{[p^e]}.$

\section{Regularity and regular sequences}

Next result can be found in \cite[16.3.9]{BroSha}.

\begin{prop}\label{exact sequences and regularity}
Let $R=\oplus_{n\geq 0} R_n$ be a positively graded, homogeneous commutative ring, and let $M$ be a non--zero, finitely generated graded $R$--module. Moreover, let $a_1,\ldots, a_h$ ($h\geq 1$) be homogeneous elements of $R$ with $\deg (a_i)=n_i\geq 1$ for all $1\leq i\leq h,$ and suppose that $a_1,\ldots ,a_h$ is an $M$--sequence. Then, the following assertions hold.

\begin{enumerate}[(i)]

\item One has that
\[
\reg^1\left(M/(a_1,\ldots, a_h)M\right)\leq\reg^1 (M)+\left(\sum_{i= 1}^h n_i\right)-h\leq\reg\left(M/(a_1,\ldots, a_h)M\right).
\]

\item Both inequalities are equalities if $\grade_M (\mathfrak{m})\geq h+1.$

\end{enumerate}
\end{prop}
\newcommand{\etalchar}[1]{$^{#1}$}
\begin{thebibliography}{HHH{\etalchar{+}}10}

\bibitem[{\`A}M20]{AlvarezMontaner2020}
J.~{\`A}lvarez~Montaner.
\newblock Local cohomology of binomial edge ideals and their generic initial
  ideals.
\newblock {\em Collect. Math.}, 71(2):331--348, 2020.

\bibitem[AMBZ20]{AlvarezBoixZarzuelaspectralsequences}
J.~\`Alvarez~Montaner, A.~F. Boix, and S.~Zarzuela.
\newblock On some local cohomology spectral sequences.
\newblock {\em Int. Math. Res. Not. IMRN}, (19):6197--6293, 2020.

\bibitem[BBR07]{BrunBrunsRomer2007}
M.~Brun, W.~Bruns, and T.~R{\"o}mer.
\newblock Cohomology of partially ordered sets and local cohomology of section
  rings.
\newblock {\em Adv. Math.}, 208(1):210--235, 2007.

\bibitem[BH98]{BrunsHerzog1993}
W.~Bruns and J.~Herzog.
\newblock {\em Cohen-{M}acaulay rings}, volume~39 of {\em Cambridge Studies in
  Advanced Mathematics}.
\newblock Cambridge University Press, Cambridge, revised edition, 1998.

\bibitem[BLS13]{BrodmannLinhSeiler2013}
M.~Brodmann, C.~H. Linh, and M.-H. Seiler.
\newblock Castelnuovo-{M}umford regularity of annihilators, {E}xt and {T}or
  modules.
\newblock In {\em Commutative algebra}, pages 207--236. Springer, New York,
  2013.

\bibitem[BS13]{BroSha}
M.~P. Brodmann and R.~Y. Sharp.
\newblock {\em Local cohomology: an algebraic introduction with geometric
  applications}, volume 136 of {\em Cambridge Studies in Advanced Mathematics}.
\newblock Cambridge University Press, Cambridge, second edition, 2013.

\bibitem[Cav07]{Cavigliatensorproducts}
G.~Caviglia.
\newblock Bounds on the {C}astelnuovo-{M}umford regularity of tensor products.
\newblock {\em Proc. Amer. Math. Soc.}, 135(7):1949--1957, 2007.

\bibitem[CHH11]{ChardinHaHoa2011}
M.~Chardin, D.~T. Ha, and L.~T. Hoa.
\newblock Castelnuovo-{M}umford regularity of {E}xt modules and homological
  degree.
\newblock {\em Trans. Amer. Math. Soc.}, 363(7):3439--3456, 2011.

\bibitem[Chr00]{Gdimensionbook}
L.~W. Christensen.
\newblock {\em Gorenstein dimensions}, volume 1747 of {\em Lecture Notes in
  Mathematics}.
\newblock Springer-Verlag, Berlin, 2000.

\bibitem[DANT02]{DivaaniAazarNaghipourTousi2002}
K.~Divaani-Aazar, R.~Naghipour, and M.~Tousi.
\newblock Cohomological dimension of certain algebraic varieties.
\newblock {\em Proc. Am. Math. Soc.}, 130(12):3537--3544, 2002.

\bibitem[DMV24]{DaoMaVarbaro}
H.~Dao, L.~Ma, and M.~Varbaro.
\newblock Regularity, singularities, and h--vector of graded algebras.
\newblock {\em Trans. Amer. Math. Soc.}, 377(3):2149--2167, 2024.

\bibitem[GS13]{M2}
Daniel~R. Grayson and Michael~E. Stillman.
\newblock Macaulay2, a software system for research in algebraic geometry.
\newblock Available at http://www.math.uiuc.edu/Macaulay2/, 2013.

\bibitem[HH06]{HoaHyry2006}
L.~T. Hoa and E.~Hyry.
\newblock Castelnuovo-{M}umford regularity of canonical and deficiency modules.
\newblock {\em J. Algebra}, 305(2):877--900, 2006.

\bibitem[HH08]{HaHoa2008}
D.~T. Ha and L.~T. Hoa.
\newblock Castelnuovo-{M}umford regularity of some modules.
\newblock {\em Comm. Algebra}, 36(3):992--1004, 2008.

\bibitem[HHH{\etalchar{+}}10]{binomialedgeideals}
J.~Herzog, T.~Hibi, F.~Hreinsd{\'o}ttir, T.~Kahle, and J.~Rauh.
\newblock Binomial edge ideals and conditional independence statements.
\newblock {\em Adv. in Appl. Math.}, 45(3):317--333, 2010.

\bibitem[HHO18]{HerzogHibiOhsugibinomialidealsbook}
J.~Herzog, T.~Hibi, and H.~Ohsugi.
\newblock {\em Binomial ideals}, volume 279 of {\em Graduate Texts in
  Mathematics}.
\newblock Springer, Cham, 2018.

\bibitem[Hol19]{HolmesgeneralizedSerrecondition}
B.~Holmes.
\newblock A generalized {S}erre's condition.
\newblock {\em Comm. Algebra}, 47(7):2689--2701, 2019.

\bibitem[IR07]{IchimRomer2007}
B.~Ichim and T.~R{\"o}mer.
\newblock On toric face rings.
\newblock {\em J. Pure Appl. Algebra}, 210(1):249--266, 2007.

\bibitem[KM11]{KumminiMurai2011}
M.~Kummini and S.~Murai.
\newblock Regularity of canonical and deficiency modules for monomial ideals.
\newblock {\em Pacific J. Math.}, 249(2):377--383, 2011.

\bibitem[MM24]{MaddoxMiller}
K.~Maddox and L.~E. Miller.
\newblock Generalized {F}--depth and graded nilpotent singularities.
\newblock {\em Comm. Algebra}, 52(11):4724--4748, 2024.

\bibitem[MS04]{MaclaganSmith2004}
D.~Maclagan and G.~G. Smith.
\newblock Multigraded {C}astelnuovo-{M}umford regularity.
\newblock {\em J. Reine Angew. Math.}, 571:179--212, 2004.

\bibitem[MS05]{MaclaganSmith2005}
D.~Maclagan and G.~G. Smith.
\newblock Uniform bounds on multigraded regularity.
\newblock {\em J. Algebraic Geom.}, 14(1):137--164, 2005.

\bibitem[MT09]{MuraiTeraihvectors}
S.~Murai and N.~Terai.
\newblock {$h$}-vectors of simplicial complexes with {S}erre's conditions.
\newblock {\em Math. Res. Lett.}, 16(6):1015--1028, 2009.

\bibitem[Ngu12a]{HopNguyenthesis}
D.~H. Nguyen.
\newblock {\em Homological and combinatorial properties of toric face rings}.
\newblock PhD thesis, Universit\"at Osnabr\"uck, 2012.
\newblock Available at \url{tinyurl.com/4vzj5e46}.

\bibitem[Ngu12b]{HopNguyen2012}
D.~H. Nguyen.
\newblock Seminormality and local cohomology of toric face rings.
\newblock {\em J. Algebra}, 371:536--553, 2012.

\bibitem[Oht11]{Ohtanibinomialedgeideals}
M.~Ohtani.
\newblock Graphs and ideals generated by some 2-minors.
\newblock {\em Comm. Algebra}, 39(3):905--917, 2011.

\bibitem[San90]{Santoni}
L.~Santoni.
\newblock Horrocks' question for monomially graded modules.
\newblock {\em Pacific J. Math.}, 141(1):105--124, 1990.

\bibitem[SZ14]{SchenzelZafarbinomialedgeideals}
P.~Schenzel and S.~Zafar.
\newblock Algebraic properties of the binomial edge ideal of a complete
  bipartite graph.
\newblock {\em An. \c{S}tiin\c{t}. Univ. ``Ovidius'' Constan\c{t}a Ser. Mat.},
  22(2):217--237, 2014.

\bibitem[Vir19]{Virilirankfunctions}
S.~Virili.
\newblock On the relation between length functions and exact {S}ylvester rank
  functions.
\newblock {\em Topol. Algebra Appl.}, 7(1):69--74, 2019.

\bibitem[Wei94]{Weibel1994}
C.~A. Weibel.
\newblock {\em An introduction to homological algebra}, volume~38 of {\em
  Cambridge Studies in Advanced Mathematics}.
\newblock Cambridge University Press, Cambridge, 1994.

\end{thebibliography}
